\documentclass[12pt]{article}
\usepackage[T1]{fontenc}
\usepackage[utf8]{inputenc}
\usepackage{lmodern}
\usepackage{textcomp}
\usepackage{enumitem}
\usepackage{geometry}
\geometry{margin=1in}
\begin{document}
\begin{center}
Answer Key - Economics - Undergraduate Level Assessment 5 \
\small (Answers are ordered exactly as in the question paper.)
\end{center}

\section*{Multiple Choice}
\begin{enumerate}[label=\arabic*.]
\item \textbf{Answer:} A
\\ \textit{Options:} A) Higher taxes on household income; B) Increased government spending; C) Increased taxes on corporate profits; D) Increased import tariffs; E) Decreased interest rates
\\ \textit{Note:} Higher taxes on household income are a supply-side fiscal policy because they can incentivize increased savings and investment by reducing disposable income, thereby aiming to enhance overall economic productivity and growth.
\item \textbf{Answer:} A
\\ \textit{Options:} A) MC that first falls, but eventually rises, as output increases.; B) AP that first falls, but eventually rises, as output increases.; C) TC that first falls, but eventually rises, as output increases.; D) AVC that first rises, but eventually falls, as output increases.; E) MP that first falls, but eventually rises, as output increases.
\\ \textit{Note:} The Law of Diminishing Marginal Returns states that as additional units of a variable input are added to a fixed input, the additional output produced from each new unit will eventually decrease, leading to a marginal cost (MC) that initially falls due to increased efficiency but ultimately rises as the fixed input becomes a limiting factor.
\item \textbf{Answer:} A
\\ \textit{Options:} A) Both the price level and real GDP fall.; B) The price level rises but the change in real GDP is uncertain.; C) The price level falls but the change in real GDP is uncertain.; D) Real GDP falls but the change in the price level is uncertain.; E) The price level rises but real GDP falls.
\\ \textit{Note:} The correct answer is misleading because when both short-run aggregate supply and aggregate demand increase, the price level may rise or fall depending on the magnitude of the shifts, but real GDP is likely to increase due to higher output resulting from the increased demand and supply.
\item \textbf{Answer:} A
\\ \textit{Options:} A) Rising per capita GDP in China increasing imports from the United States; B) Lower per capita GDP in China decreasing imports from the United States; C) Rising unemployment rates in the United States; D) Higher price levels in the United States relative to China; E) Decrease in the American GDP
\\ \textit{Note:} Rising per capita GDP in China leads to increased consumer purchasing power and demand for imports, including U.S. goods, which boosts the demand for the U.S. dollar and causes its appreciation relative to the Chinese yuan.
\item \textbf{Answer:} A
\\ \textit{Options:} A) Overproduction and wastage of crops; B) Lack of alternative employment; C) Excessive government intervention in agriculture; D) Inefficient farming methods; E) Low income of farmers
\\ \textit{Note:} The correct answer highlights that overproduction and wastage of crops occur when resources are not allocated efficiently in agriculture, leading to a surplus that surpasses market demand and results in economic inefficiencies.
\end{enumerate}

\section*{True / False}
\begin{enumerate}[label=\arabic*.]
\setcounter{enumi}{5}
\item \textbf{Answer:} False
\\ \textit{Options:} A) True; B) False
\\ \textit{Note:} Cyclical unemployment is primarily caused by fluctuations in the business cycle, such as recessions and economic downturns, rather than by technological advancements that lead to job obsolescence.
\item \textbf{Answer:} False
\\ \textit{Options:} A) True; B) False
\\ \textit{Note:} Labor is often considered more expensive than land and capital in the United States due to higher wage levels and associated labor costs, such as benefits and regulations.
\item \textbf{Answer:} True
\\ \textit{Options:} A) True; B) False
\\ \textit{Note:} The statement is true because with a marginal propensity to consume (MPC) of 0.75, the multiplier effect is 4 (1/(1-MPC)), meaning that an initial investment of $75 million will lead to a total income change of $300 million (75 million x 4).
\item \textbf{Answer:} False
\\ \textit{Options:} A) True; B) False
\\ \textit{Note:} The statement is false because the average total cost (ATC) curve can be above, equal to, or below the marginal cost (MC) curve, depending on the level of production, particularly when the MC is less than the ATC during the increasing returns to scale phase.
\item \textbf{Answer:} False
\\ \textit{Options:} A) True; B) False
\\ \textit{Note:} The statement is false because total wealth is not solely determined by annual income and interest rate; it also depends on the accumulation of assets and liabilities, which are not provided in the statement.
\end{enumerate}

\section*{Long Form Response}
\begin{enumerate}[label=\arabic*.]
\setcounter{enumi}{10}
\item \textbf{Answer:} In a perfectly competitive market operating below the minimum average total cost in the long run, firms are unable to cover their total costs, leading to economic losses. As a result, some firms will begin to exit the industry, reducing the overall supply in the market. This decrease in supply creates upward pressure on prices, which will eventually rise towards the breakeven point, where firms can cover their costs. In the long run, the market will reach a new equilibrium at a price that allows remaining firms to operate at a sustainable level. Thus, firm behavior in response to losses and market adjustments ultimately drives the price back to a point that supports viability in the industry.
\\ \textit{Note:} correct because it accurately describes how firms in a perfectly competitive market suffering economic losses will exit the industry, leading to a decrease in supply that drives prices upward until they reach a level where firms can cover their costs and achieve long-run equilibrium.
\item \textbf{Answer:} Permanent inflation has significant consequences on interest rates, primarily through the mechanism of anticipated inflation. When inflation is expected to persist, lenders and investors adjust the nominal interest rates they charge to account for the erosion of purchasing power over time. This means that the nominal interest rates become effectively discounted by the anticipated rate of inflation, leading to real interest rates that may remain lower than they would be in a stable price environment. Consequently, this can distort saving and investment decisions, as individuals and businesses may seek to minimize their exposure to inflation risk, potentially leading to decreased economic growth. Thus, understanding the relationship between inflation expectations and nominal interest rates is crucial for both policymakers and economic agents in navigating an inflationary landscape.
\\ \textit{Note:} The answer correctly identifies that permanent inflation leads to an adjustment in nominal interest rates to compensate for anticipated inflation, which ultimately influences real interest rates and reflects the erosion of purchasing power in the economy.
\end{enumerate}

\vfill
\noindent\textit{End of Answer Key}
\end{document}